% 本文件由 LaTeX 排版 Agent 根据有效 translation.json 自动生成。
% author=Jerome; work_id=3002; title=序言

\chapter{序言}

% structure: p0001 source=h1 target=omitted reason=page-title-already-rendered-by-parent

% structure: p0002 source=h2 target=section reason=relative-heading-rank; base=h2

\section{尤西比乌斯\WorkTitle{编年史}序言}

\Quote{\WorkTitle{编年史}}是一部普世史著作，记载了从\Person{亚巴郎}蒙召和奥林比亚德开始的年代。关于它的说明，读者可参阅\Person{萨蒙}博士在\Quote{\WorkTitle{基督教古物词典}}中的文章。该书由\Person{热罗尼莫}于381-82年在君士坦丁堡翻译，当时他正为参加大公会议而逗留于此。这篇序言表明，\Person{热罗尼莫}已经开始意识到旧约各种译本所带来的困难，以及回到希伯来文原文的必要性。\par

\Person{热罗尼莫}致其友人\Person{文森提乌斯}和\Person{加里埃努斯}，问安：\par

1. 学者们长久以来习惯于通过将希腊作家的作品译成拉丁文来锻炼心智，更困难的是，还要在格律的进一步限制下翻译著名诗人的诗篇。因此，我们的\Person{图利乌斯}逐字翻译了\Person{柏拉图}的整卷书；在出版了一部用六音步诗体写成的\Person{阿拉托斯}（如今可被视为罗马人）的版本后，他又以\Person{色诺芬}的\WorkTitle{家政论}自娱。在这后一部作品中，那金色的雄辩之河一再遇到障碍，河水在其周围激荡飞溅，以至于不熟悉原文的人不会相信他们读的是\Person{西塞罗}的文字。这不足为奇！要追随他人的思路处处守住界限是困难的。要在译作中保持优雅与韵味而不受损伤是一项艰巨的任务。某个词有力地表达了一个思想；我却没有自己的词来传达意思；当我在寻求满足原意时，可能绕了很远的路却只完成了行程的一小段。我们还需考虑词序的变换、格的差异、修辞格的多样性，以及最后，那种独特的、可谓语言本身的地道习语。直译听起来荒谬；另一方面，如果我被迫改变顺序或词语本身，就会显得我放弃了译者的职责。\par

2. 因此，我亲爱的\Person{文森提乌斯}，还有你，\Person{加里埃努斯}，我所爱如同自己灵魂的人，我恳求你们，无论这份仓促之作价值如何，请以朋友而非批评家的心情来阅读。我之所以如此恳切地请求，是因为如你们所知，我极为迅速地口述给我的书记员；而圣经的见证表明这项任务何其困难；因为七十子希腊译本并未保留原有的风味。正是这一点促使了\Person{阿奎拉}、\Person{辛马库斯}和\Person{提奥多提翁}的努力；他们工作的结果是为同一部作品赋予了截然不同的特点：一人力求逐字对应，另一人力求总体意思，而第三人希望避免与古人任何重大的差异。第五、第六和第七种版本，虽然无人知晓其作者是谁，却展现出一种如此令人愉悦的多样性，以至于尽管匿名，却赢得了权威的地位。因此，有些人的看法甚至发展到认为圣经听起来有些粗涩刺耳；这是因为我所提到的那些人不知道这些作品是从希伯来文翻译过来的，因此，他们只看表面而不看实质，在发现语言所装饰的优美身体之前，便对那简陋的外表感到战栗。事实上，有什么比\BibleBook{圣咏}{集}更富有韵律呢？它像我们自己的\Person{弗拉库斯}和希腊的\Person{品达}的作品一样，时而以抑扬格轻快前进，时而以响亮的阿尔凯奥斯格律流淌，时而升华为萨福格律，时而以半音步行进。有什么比\BibleBook{申命}{纪}和\BibleBook{依撒意亚}{}的旋律更可爱？有什么比\Person{撒罗满}的话语更庄重？有什么比\BibleBook{约伯}{传}更完美？所有这些，正如\Person{若瑟夫}和\Person{奥利金}告诉我们的，都是以六音步和五音步诗体写成的，并在他们自己的民族中流传。当我们读希腊文时，它们具有\Emph{某些}意义；而读拉丁文时则全然不连贯。但如果有人认为语言的美不会因翻译而受损，就请他把\Person{荷马}逐字译成拉丁文。我甚至要说，如果他把这位作者用他自己的语言译成散文，词序会显得可笑，而这位最雄辩的诗人几乎成为哑巴。\par

3. 所有这一切的用意何在？我不希望你们觉得奇怪，如果我们在这里或那里绊倒；如果语言迟缓；如果辅音丛生或元音出现缺口；或者因叙事的紧凑而显得局促。最博学的人也在这同一任务上辛劳过；而且，除了所有人都会经历且我们已指出伴随所有翻译工作的困难之外，不可忘记，我们还面临一个特殊的困难，因为历史是多方面的，充满了拉丁人不知道的野蛮名称、情况，以及错综复杂的日期，评注标记与事件和数字混杂在一起，以至于几乎分辨词序比认识所述之事更难。\par

\EditorialNote{以下是一段长文，展示了按照某种排列方式，日期通过特定颜色区分，以表示属于所涉及的某个王国历史。这段文字在没有彩色图表的情况下似乎难以理解，并且除非在研究带有原始排列的书，否则没有用处。}\par

我很清楚，会有许多人，以他们惯于普遍诋毁的爱好（摆脱此事的唯一办法就是什么都不写），将他们的毒牙咬入这部书。他们会挑剔日期，改变顺序，驳斥事件的准确性，筛分音节，并且，正如常常发生的那样，会把抄写员的疏忽归咎于作者。我若反驳他们，说他们可以不必阅读，如果他们选择不读，这本是我的权利；但我宁愿让他们心平气和地离开，好让他们将应得的荣誉归于希腊作者，并认识到我们所负责的任何增补都是取自其他极有声望的人。事实是，我部分履行了翻译者的职责，部分履行了作者的职责。我极其忠实地翻译了希腊文部分，同时又添加了一些在我看来已被忽略的内容，尤其是在\Place{罗马}史方面，这部书的作者\OriginalTerm{欧瑟比}{Eusebius}在我看来只是略略触及；这与其说是由于无知——因为他是有学问的人——不如说是由于他使用希腊文写作，认为这些内容对他的同胞无关紧要。同样，从\OriginalTerm{尼努斯}{Ninus}和\OriginalTerm{亚巴郎}{Abraham}一直到\OriginalTerm{特洛伊}{Troy}被俘，译文仅来自希腊文。从特洛伊到\OriginalTerm{君士坦丁}{Constantine}二十年，有许多内容，时而单独添加，时而混合其中，是我从\OriginalTerm{特朗基鲁斯}{Tranquillus}和其他著名史学家那里极为勤奋地搜集而来的。此外，从上述君士坦丁那一年到\OriginalTerm{瓦伦斯}{Valens}皇帝第六次执政官任期和\OriginalTerm{瓦伦提尼安}{Valentinianus}第二次执政官任期之间的部分，完全是我自己的作品。我满足于在此结束，将剩余的时期，即\OriginalTerm{格拉提安}{Gratianus}和\OriginalTerm{德奥多西}{Theodosius}时期，保留给更广阔的历史综述；这并不是因为我害怕自由而真实地谈论活着的人，因为敬畏天主就驱除对人的畏惧；而是因为当我们的国家仍暴露在野蛮人的狂怒之下时，一切都在混乱之中。\par

% structure: p0010 source=h2 target=section reason=relative-heading-rank; base=h2

\section{译\Person{奥力振}\WorkTitle{雅歌}两篇讲道辞序}

写于\Place{罗马}，公元383年。\par

\Person{热罗尼莫}致至圣教宗\Person{达玛稣}：\par

\Person{奥力振}在其其他著作中已超越所有人，而在\WorkTitle{雅歌}中则超越了自己。他为此写了十卷，约有两万行，在这些卷中，他首先讨论了七十贤士译本，然后讨论了\Person{阿奎拉}、\Person{息玛古}和\Person{德敖多齐}的译本，最后讨论了他声称在\Place{亚特里}海岸发现的第五个译本，其文采之壮丽与充实，在我看来他实现了诗中所说的：\Quote{君王带我进了他的内室。}我将那部著作搁置一边，因为要将如此伟大的著作译成拉丁文，即使能够得体地完成，也需要几乎无限的闲暇、劳力和金钱；而我翻译了这两篇短论，是他以每日演讲的形式为那些仍像婴孩和吃奶者的人创作的，我注重忠实而非优雅。当较小的作品已令您如此满意时，您可以想见较大的作品具有何等巨大的价值。\par

% structure: p0014 source=h2 target=section reason=relative-heading-rank; base=h2

\section{\WorkTitle{希伯来名字}序}

本书的起源和范围在序言本身中已有描述。它写于388年，即\Person{热罗尼莫}定居\Place{白冷}两年后。他在抵达\Place{巴勒斯坦}后立即着手（三年前）提高他的希伯来语知识，以准备\WorkTitle{旧约}的翻译，该翻译始于391年。因此，本书以及随后的两部作品，可视为为拉丁通行本所做的预备性研究的记录。\par

\Person{斐罗}，犹太人中最博学的人，据\Person{奥力振}称，他做了我现在正在做的事；他撰写了一部\WorkTitle{希伯来名字}的书，按其首字母分类，并将每个词源列于一旁。我原拟将这部著作译成拉丁文。它在希腊世界很著名，在所有图书馆都能找到。但我发现各抄本彼此极不一致，顺序混乱，因此我判断最好什么也不说，而不是写出会受公正谴责的东西。然而，这类作品看起来可能有用；我的朋友\Person{路普利亚诺}和\Person{瓦莱利亚诺}敦促我尝试，因为他们认为我在希伯来语知识方面已有一些进步。因此，我按顺序通过了全部\WorkTitle{圣经}各卷，在我现在对古老架构的复原中，我认为我创作了一部可能对希腊人和拉丁人都有价值的作品。\par

我在此序言中请求读者注意，如果他在本作品中发现有遗漏之处，那是保留在另一作品中提及的。我目前手头有一本\WorkTitle{希伯来问题}的书，这是一项全新的工作，迄今为止在希腊人或拉丁人中从未听说过。我说这话，并非狂妄地抬高自己的作品，而是因为我知道我在其中投入了多少劳动，并希望激发那些知识不足的人去阅读它。我建议所有希望拥有那部作品和本作品，以及我将要出版的\WorkTitle{希伯来地名}一书的人，不要理会犹太人及其所有的恼怒爆发。此外，我还添加了\WorkTitle{新约}中词语和名字的含义，以使建筑得到最后的润色并完好无缺。我也想在这件事上效法\Person{奥力振}，除了无知者外，所有人都承认他是仅次于宗徒的教会最伟大的导师；因为在这部作品（它属于他天才最高贵的纪念碑之一）中，他作为一名基督徒，努力补充了身为犹太人的\Person{斐罗}所遗漏的内容。\par

% structure: p0018 source=h2 target=section reason=relative-heading-rank; base=h2

\section{\WorkTitle{希伯来地名位置与名称}序}

\EditorialNote{写于公元388年。}\par

\Person{欧瑟伯}，从他以真福殉道者\Person{庞斐禄}的名字取得了第二个名字，在撰写了他的十卷\WorkTitle{教会史}、\WorkTitle{年代纪}（我曾出版过其拉丁文译本）、详述各族名称以及犹太人古时和现今所称呼的名称的书、\Place{犹太地区}的地志和各支派分得的土地，连同\Place{耶路撒冷}及其圣殿的地图（并附有极简短的说明）之后，在生命终了时，为这部小作品付出了辛劳，其设计是从圣经中为我们收集几乎所有城市、山岳、河流、村庄及其他地点的名称，不论它们是否保留原名，或后来有所改变，或在一定程度上被讹误。我接过了这位可钦之人的作品，并翻译了它，遵循希腊人的编排，按字母顺序取词，但删除了那些似乎不值得提及的名字，并作了相当多的改动。我在我的\WorkTitle{年代纪}翻译的序言中已一次性解释了我的方法，在那里我说我可以同时被称为译者和一部新作品的作者；但我特别重复这一点，因为一个几乎没有丝毫文字素养的人竟敢冒险将这本书翻译为拉丁文，不过他的语言几乎不能称为拉丁文。细心的读者但凡将其与我译本比较，便会看出其学问的匮乏。我并不自诩有凌驾云霄的风格；但我希望能超越那匍匐于地的风格。\par

% structure: p0021 source=h2 target=section reason=relative-heading-rank; base=h2

\section{\WorkTitle{希伯来问题书}序言}

\EditorialNote{写于公元三八八年。}\par

书序的目的是阐述后续作品的主旨；但此刻我不得不先回答针对我提出的反对意见。我的情况有点像\Person{泰伦提乌斯}，他把剧本的剧场开场白变成了自我辩护。我们有一个\Person{路斯基乌斯·拉努维努斯}，就像那个折磨他的人一样，指控这位诗人如同掠夺国库的强盗。\Place{曼图阿}的诗人也遭受了同样的情形；他极为精确地翻译了几行\Person{荷马}史诗，他们就说他不过是古人的剽窃者。但他回答说，从\Person{赫拉克勒斯}手中夺过他的棍棒，绝非力量小的证明。为何，就连\Person{图利乌斯}——这位罗马修辞学的巅峰、演说家之王和拉丁语的荣耀——也被希腊人控告贪污。因此，若像我这样一个可怜的小人物暴露于那些践踏我们珍珠的卑劣猪猡的哼哼声，我并不感到惊讶，因为一些最博学的人——他们的声望本应平息恶意之声——也曾感受到嫉妒的火焰。诚然，这发生得颇有几分公道，对那些其雄辩充满剧院和元老院、公民大会和讲坛的人来说；鲁莽总是招惹诽谤，并且（正如\Person{贺拉斯}所说）：\par

最高的山峰招致\newline{}闪电的打击。\par

但我身处一隅，远离城市和市场，远离拥挤法庭的争吵；然而，即便如此（正如\Person{昆体良}所说），恶意还是找到了我。因此，我恳求读者，\par

若有什么人，若有一个人，\newline{}被强烈的渴望所陶醉，将要阅读这些文字，\par

不要期待在这些\WorkTitle{希伯来问题书}中有雄辩或演说家的优美——我打算就所有圣经书卷撰写此书；而是希望他亲自替我回答那些诋毁我的人，并告诉他们，一种新体裁的作品是可以要求一些宽容的。我贫穷且地位低微；我既无财富，也不认为若有人馈赠我就应该接受；同样，让我告诉他们，他们不可能同时拥有基督的财富，即圣经的知识，和世俗的财富。因此，我的简单目标首先是指出那些怀疑希伯来圣经有误之人的错误，其次，通过参照原始权威，纠正那些显然充斥于希腊文和拉丁文抄本的错误；此外，当从拉丁词语的发音无法明了时，通过用俗语进行意译，来解释事物、名字和国度的词源。为使学生更容易注意到这些修正，我首先计划如我现在所能做的那样，列出正确的读法本身，然后通过将后来的读法与它比较，指出哪些被省略、添加或改变了。我并非如龇牙咧嘴的恶意所声称的那样，要定罪于七十贤士译本的错误，也不认为我的劳动是对他们工作的贬低。事实是，既然他们的工作是为\Person{亚历山大}的\Person{托勒密}王所做，他们便不愿揭示圣经所蕴含的一切奥秘，尤其是那些预告基督来临的奥秘，唯恐那位因犹太人被认为崇拜独一天主而尊重他们的人，会以为他们崇拜第二位神。但我们发现，福音作者，甚至我们的主救主，以及宗徒\Person{保禄}，也引用了许多出自旧约的经文，这些经文并不包含在我们的抄本中；关于这些，我将在适当的地方更充分地阐述。但由此显然可见，那些最符合新约权威言辞的手抄本是最好的。此外，\Person{若瑟夫}叙述了七十贤士译者的故事，但报道他们只翻译了\Person{梅瑟}五书；我们也承认这些比其余部分更符合希伯来文。而且，后来出现的译者——我指\Person{阿基拉}、\Person{辛马库斯}和\Person{狄奥多田}——的译本与我们使用的译本大不相同。\par

我只有一句话要说，这或许能使我的反对者平静下来。外国货只应输入有需求的地方。乡下人没有义务购买香脂、胡椒和椰枣。至于\Person{奥力振}，我什么也不说。他的名字（若我可拿小事与大事相比）比我的名字更招人忌恨，因为他在面向普通群众的讲道中虽遵循通行译本，但在他的\WorkTitle{卷册}中——即他对圣经更详尽的讨论中——却以希伯来文为真理，尽管有自己阵营的拥护，却偶尔求助外邦语言作为盟友。我只想就此说一句：我宁愿拥有他对圣经的知识，纵使伴随他名下的一切忌恨，也毫不在乎这些阴影和幽灵般的鬼怪——据说它们的本性是在黑暗角落里叽叽喳喳，吓唬婴儿。\par

% structure: p0029 source=h2 target=section reason=relative-heading-rank; base=h2

\section{训道篇注释序}

\Emph{\Person{保拉}与\Person{欧斯托基姆}，\Place{白冷}，公元三八八年。}\par

我记得，大约五年前，当我仍住在\Place{罗马}时，我曾向圣善的\Person{布雷西拉}诵读\BibleBook{训道}{篇}，以激发她轻看这尘世的一切，视她在世上所见为虚无；她曾请求我将对其中所有晦涩段落的评述写成一篇短注，以便我缺席时，她仍能理解所读。然而，她猝然离我们而去，正当为我们的事业束腰预备时；我们不被视为配得有这样一位人生伴侣；因此，\Person{保拉}和\Person{欧斯托基姆}，我在如此创伤之下保持沉默。但现在，我生活在较小的\Place{白冷}团体中，我将对亡者的纪念以及对你们的欠债偿还。我只想指出这一点：我没有跟随任何人的权威。我直接译自希伯来文，尽可能使我的措辞接近七十贤士译本的形式，但仅在其与希伯来文相差不远的段落中。我也偶尔参考了\Person{阿基拉}、\Person{息玛古}和\Person{狄奥多田}的译本，但既不致因过多新奇而惊扰勤勉的学生，也不让我的注解追随意见的支流，违背我良心的确信，偏离真理的源头。\par

% structure: p0032 source=h2 target=omitted reason=duplicate-page-title

此译本于公元三八二年至三八五年间在\Place{罗马}完成。仅存的序言是为福音译本所作，但\Person{热罗尼莫}也曾引用和谈论他对其他部分的译本。此项工作是在\Person{达玛苏教宗}的请求和核准下进行的，教宗曾于公元三八三年就圣经评注的某些要点咨询过\Person{热罗尼莫}，显然在同一年敦促他借助希腊原文修订当时通行的拉丁译本。应当指出，\Person{热罗尼莫}的目的是\Quote{修订旧拉丁译本，而非制作新译本}。当\Person{奥斯定}向他表达对\InnerQuote{他的\Emph{翻译}福音}的感谢时，他默然纠正，以\InnerQuote{新约的\Emph{校订}}替代了这一说法。然而，尽管他为自己设定了这个有限的目标，正如他所描述的，所引入的各种讹误形式如此众多，以致旧译本与修订本（\Person{热罗尼莫}本）之间的差异始终清晰而显著。参看\Person{韦斯科特}在\WorkTitle{圣经辞典}中关于武加大译本的条目，以及\Person{弗里曼特尔}在\WorkTitle{基督教传记辞典}中关于\Person{热罗尼莫}的条目。\par

% structure: p0034 source=h2 target=section reason=relative-heading-rank; base=h2

\section{四福音}

\Emph{\Person{达玛苏教宗}，公元三八三年。}\par

你催促我修订旧拉丁文译本，并如同审判那些如今散布于全世界的圣经抄本一般；且因它们彼此互异，你要我断定哪些与希腊原文相符。这项工作出于真爱，但同时又危险而冒昧；因为在判断他人时，我必须甘心受所有人判断；我怎敢在世界的语言已至垂暮之年时改变它，并使之重返其初创之时呢？有哪一个有学问或没有学问的人，当他拿起这卷书，发觉他所读的不符合他固定的口味时，不会立即破口大骂，称我为伪造者和亵渎者，因为我竟敢在古书里添加任何东西，或作任何改动或修正？如今有两个安慰的思考使我能够忍受这憎恶——首先，这命令是由您，至高主教所赐；其次，即使按那些辱骂我们之人的说法，与早期抄本不符的读法也不能是正确的。因为我们若信靠拉丁文本，就应由我们的对手告诉我们\Emph{哪一个；}因为几乎有多少抄本就有多少文本形式。另一方面，我们若从\Emph{许多}文本的比较中收集真理，为何不回到希腊原文，修正由不准确的译者引入的错误，以及自信而无知的批评者的拙劣窜改，更不用说那些半睡半醒的抄写者所插入或更改的一切呢？我不是在讨论旧约，它由七十贤士译成希腊文，经过三步的传承到达我们。我不问\Person{阿奎拉}和\Person{息马古}所信，或为何\Person{特敖多田}在古人和今人之间采取中间路线。我愿将那得到宗徒认可的译本视为真译本。我现在说的是新约。新约无疑是用希腊文写成的，但\Person{玛窦}宗徒的作品除外，他是第一个将基督的福音写下来，并在\Place{犹太地}用希伯来文字发表其作品的人。我们必须承认，在我们语言中的译本充满了分歧，如今河流分成不同的水道，我们必须回到源头。我略过那些与\Person{路济安}和\Person{赫西丘}的名字相关的抄本，它们的权威被一小撮好争辩的人固执地维护。显然，在七十贤士的工作之后，这些作者无法修正旧约中的任何东西；而矫正新约也是无用的，因为已经存在于许多民族语言中的圣经译本显示他们的添加是虚假的。因此，在这简短的序言中，我只承诺四福音书，它们应取以下顺序：\Person{玛窦}、\Person{玛尔谷}、\Person{路加}、\Person{若望}，正如它们通过比较希腊抄本而被修订一样。只使用了早期的抄本。但为了避免与我们习惯阅读的拉丁文有太大的分歧，我用了较为克制的笔触，只修正了那些似乎传达不同含义的段落，其余的我保留原样。\par

序言以描述由\Person{欧瑟伯}编制并由\Person{热罗尼莫}翻译的词语列表作为结尾，旨在显示哪些段落出现在两卷或多卷福音书中。\par

\SourceRecord{Translated by W.H. Fremantle, G. Lewis and W.G. Martley.；Nicene and Post-Nicene Fathers, Second Series；Vol. 6.；Edited by Philip Schaff and Henry Wace.；Buffalo, NY: Christian Literature Publishing Co.,；1893.；<http://www.newadvent.org/fathers/3002.htm>.}
